% Generated draft: requires mathematical and source review.
\documentclass[11pt]{article}
\usepackage[margin=25mm]{geometry}
\usepackage{fontspec}
\setmainfont{Latin Modern Roman}
\usepackage{amsmath,amssymb,mathtools,unicode-math}
\setmathfont{Latin Modern Math}
\usepackage{enumitem,needspace,xurl,hyperref}
\setlength{\parindent}{0pt}
\setlength{\parskip}{5pt}
\setlength{\emergencystretch}{3em}
\newcommand{\N}{\mathbb{N}}
\newcommand{\Z}{\mathbb{Z}}
\newcommand{\Q}{\mathbb{Q}}
\newcommand{\R}{\mathbb{R}}
\newcommand{\C}{\mathbb{C}}
\newcommand{\F}{\mathbb{F}}
\newcommand{\A}{\mathbb{A}}
\newcommand{\PP}{\mathbb{P}}
\newcommand{\E}{\mathbb{E}}
\newcommand{\eps}{\varepsilon}
\newcommand{\dd}{\,\mathrm{d}}
\begin{document}
\section*{UnsolvedMath ID: 800072}
\textbf{Source problem code:} catalogue entry
\subsection{Definitions and mathematical statement}
Let \(M\) be a four-dimensional manifold endowed with a foliation \(\mathcal F\) by two-dimensional leaves. In the smooth setting needed for the following definitions, its tangent distribution \(T\mathcal F\subset TM\) is a smooth rank-two subbundle satisfying \([\Gamma(T\mathcal F),\Gamma(T\mathcal F)]\subseteq\Gamma(T\mathcal F)\).

A smooth rank-two distribution \(\mathcal D\subset TM\) is an \emph{Engel structure} if its derived distributions
\[
\mathcal D^{2}=\mathcal D+[\mathcal D,\mathcal D],\qquad
\mathcal D^{3}=\mathcal D^{2}+[\mathcal D,\mathcal D^{2}]
\]
have ranks three and four, respectively, at every point of \(M\). Here brackets of distributions denote the pointwise spans of Lie brackets of their smooth local sections; in particular, \(\mathcal D^{3}=TM\).

\textbf{Question (Nariman--Yazdi, Question 13).} Given a four-manifold \(M\) with a foliation \(\mathcal F\) by surfaces, can \(\mathcal F\), interpreted through its tangent distribution \(T\mathcal F\), be perturbed to an Engel structure?

\textbf{Unresolved meaning of perturbation.} The supplied statement does not specify a topology, a size of perturbation, or whether a deformation is required. For example, \emph{if} perturbation means arbitrarily close approximation in the compact-open \(C^{\infty}\) topology on the space \(\Gamma(\operatorname{Gr}_{2}(TM))\) of smooth rank-two distributions, the question would ask, for every admissible pair \((M,\mathcal F)\), whether
\[
\forall\,U\text{ a neighborhood of }T\mathcal F,
\quad \exists\,\mathcal D\in U
\quad\text{such that }\mathcal D\text{ is Engel}.
\]
This is a conditional interpretation, not a uniquely determined formulation of the supplied question. The source also leaves smoothness conventions, compactness, boundary conditions, and any orientation hypotheses unspecified; no such additional restrictions are asserted here.
\subsection{Short English statement}
The question asks whether a foliation of a four-manifold by surfaces can be perturbed so that its tangent planes form an Engel distribution, whose successive Lie brackets span three and then four dimensions everywhere. The meaning and size of the perturbation, and the intended manifold and regularity hypotheses, require clarification.
\subsection{Sources}
\begin{itemize}
\item Supplied local UnsolvedMath registry record: permanent ID \(800072\), exact source problem code \texttt{catalogue entry}, title ``Foliations and Diffeomorphism Groups --- Question 13.'' The record attributes the problem to Nariman and Yazdi, with proposed year 2025.
\item Nariman and Yazdi, \emph{Problem List on Foliations and Diffeomorphism Groups} (2025), Question 13. Supplied source URL: \url{https://arxiv.org/abs/2503.18274}. The supplied provenance records access on 2026-10-16 and extraction by \texttt{arxiv\_latex}; these are supplied provenance claims, not independently verified access or source verification.
\item The supplied registry labels the problem \emph{open}. This is a supplied status label, not an independently verified mathematical or literature assessment. The source introduces the question as part of a discussion of interactions between Engel structures and other geometric structures and mentions two questions; only the surface-foliation question is supplied here.
\end{itemize}
\end{document}
